% Source: source_md/intelligence_part2_appendices_ja.md
\section{主要定理の証明スケッチ集}


本補遺は、本論文で提示した主要な定理・系・主張の証明スケッチを整理する。本論文は形式的証明体系を組まないが、各命題が論理的に成立する理由を構造的に示す。

\bigskip


\subsection{定理 3.1(非閉包定理 / Non-Closure Theorem)}


\textbf{命題}: 前提 (P1)(P2) を満たす任意の推論演算子 $F$ について、内部不動点 $x^*$ は存在しない。

\[
\forall x \in \mathcal{X}, \quad F(x) \neq x
\]

\textbf{前提の再掲}:
\begin{itemize}
\item (P1) 正当化拡張性: $F$ は $x_n$ をより正当化された $x_{n+1}$ へ写す
\item (P2) 正当化測度の存在: 測度 $J: \mathcal{X} \to \mathbb{R}$ が存在し、$J(F(x)) > J(x)$ が成立
\end{itemize}


\textbf{証明スケッチ}:

$F(x^*) = x^*$ を満たす不動点 $x^* \in \mathcal{X}$ が存在すると仮定する。

前提 (P2) より任意の $x$ について $J(F(x)) > J(x)$。

特に $x = x^*$ とすると: $J(F(x^*)) > J(x^*)$

しかし不動点の仮定より $F(x^*) = x^*$ ゆえ: $J(x^*) > J(x^*)$

これは矛盾。ゆえに不動点 $x^*$ は存在しない。$\blacksquare$

\textbf{注意}: 証明は正当化測度の単調性のみに依拠する。$F$ の具体的実装(演繹的・確率的・統計的)は問わない。

\bigskip


\subsection{系 3.2(内部停止述語の不可能性)}


\textbf{命題}: 推論内部に停止述語 $S: \mathcal{X} \to \{0, 1\}$ を定義することはできない。

\textbf{証明スケッチ}:

仮に内部停止述語 $S$ を定義するとする。$S(x) = 1$ となる状態 $x$ は、これ以上推論を進める必要がない、すなわち $F$ の作用を要しない状態でなければならない。これは $F(x) = x$ を意味し、定理 3.1 に矛盾する。

$S$ を「停止判断のための補助推論 $F'$」として定義しようとする場合、$F'$ 自体が推論演算子であり、定理 3.1 が適用される。$F'$ にも内部不動点はなく、$S$ を確定するためにはさらに $S'$ が必要となる。この連鎖は終わらない。$\blacksquare$

\bigskip


\subsection{系 3.3($C_{max}$ 不到達性との同型)}


\textbf{命題}: 内部不動点不在は VED の $C_{max}$ 不到達性と構造的に同型である。

\textbf{スケッチ}:

VED 基本式の最終項 $-\kappa_B/(C_{max} - C)$ は $C \to C_{max}$ で発散し、完全閉包への到達を動力学的に禁じる。

知性層において:
\begin{itemize}
\item 内部不動点 $F(x^*) = x^*$ ← → 正当化の完全閉包(これ以上正当化を要しない状態)
\item 定理 3.1 による内部不動点不在 ← → $C_{max}$ 不到達性
\end{itemize}


対応関係:

\begin{longtable}{@{}p{0.30\linewidth}p{0.30\linewidth}@{}}
\toprule
VED 物理層 & 知性層 \\
\midrule
\endfirsthead
\toprule
VED 物理層 & 知性層 \\
\midrule
\endhead
$C \to C_{max}$ での発散 & $F$ の内部不動点不在 \\
$\kappa_B/(C_{max}-C)$ 障壁 & 正当化の無限後退(系 3.2) \\
準閉包 $\Omega_\theta$ & 操作可能推論領域 \\
\bottomrule
\end{longtable}


両者は別領域の現象だが、同一の動力学的命題の層別表現である。$\blacksquare$

\bigskip


\subsection{系 3.4(個体閉包不可能性)}


\textbf{命題}: 個体システムは推論層 $L_F$ 単独で推論を閉じることができない。

\textbf{スケッチ}:

定理 3.1 により $L_F$ に内部不動点はない。系 3.2 により内部停止述語も定義できない。推論の停止が $L_F$ の内部からは来えない。

したがって停止は $L_F$ の外側から来なければならない。その外側の構造が $L_E$ である(\S{}5)。$L_F$ 単独では停止不能 ← → 「個体内で閉じない」。ただし $L_R$ と $L_E$ を含む三層全体では外部依存的に閉じうる(\S{}5)。$\blacksquare$

\bigskip


\subsection{定義 1.1(VED 的創発)の三例検証}


\textbf{命題}: $\mathcal{U} = \Phi(\mathcal{L}\text{の運動}, H_{ij})$ という創発の形式が三例すべてで成立する。

\textbf{スケッチ}:

\textbf{(1) Vol.2: 標準模型構造の創発}

$\mathcal{L}$: 因果ログ $C_{ij}$ と内部ベクトル $\vec{e}_i$ の運動。

$H_{ij}$: 閉包条件 $\sum \vec{e} \approx 0$、閉包エネルギー汎関数 $E(A)$。

$\mathcal{U}$: 三角形閉包(最小自己参照的 enclosure)および粒子(= $E(A)$ の局所極小)。

$\mathcal{U}$ は $\mathcal{L}$ から演繹できないが、閉包条件という $H_{ij}$ の下で $\mathcal{L}$ の運動が到達する安定状態として生じる。✓

\textbf{(2) Bio-IFGT: 生命の創発}

$\mathcal{L}$: 情報流 $J$ と状態発展 $\dot{x} = F(x;\theta)$ の運動。

$H_{ij}$: DNA 由来パラメータ $\theta$、環境との非平衡物質・エネルギー交換。

$\mathcal{U}$: 生命アトラクター $\Omega_{\text{life}}$。

化学反応の運動から、$\theta$ という拘束パラメータの下で生命という準閉包が創発する。✓

\textbf{(3) 本論文: 創発停止層 $L_E$}

$\mathcal{L}$: ログ層 $L_R$ の再帰運動。

$H_{ij}$: 観測者-環境相互拘束。

$\mathcal{U}$: 創発停止層 $L_E$(停止アトラクター)。

$L_R$ の運動は $L_E$ を論理的には含意しないが、$H_{ij}$ の下での $\mu_\theta$ による (B) 組み替えを通じて、$L_E$ が動力学的に生じる。✓

すべての例が同一の形式を満たす。$\blacksquare$

\bigskip


\subsection{主張 5.4($\Phi_{\text{ext}} \equiv \Phi_{\text{net}}$)のスケッチ}


\textbf{命題}: 外部参照場 $\Phi_{\text{ext}}$ は個体間 $L_E$ ネットワーク $\Phi_{\text{net}}$ と同型である。

\textbf{スケッチ}:

$\Phi_{\text{ext}}$ は「個体が停止を委ねる外部の共通参照基盤」として機能論的に定義される。

$\Phi_{\text{net}}$ は「複数の個体が各自の内部に持つ $L_E$ が、相互干渉によって形成するネットワーク」として構造的に定義される。

外から観察者が見ると、個体 $i$ は $\Phi_{\text{ext}}$ という単一の参照場を参照しているように見える。しかし構造的には、個体 $i$ が参照しているのは、他の個体 $j, k, \ldots$ の $L_E$ との相互干渉ネットワーク $\Phi_{\text{net}}$ である。

この機能的外観と構造的実体の一致が同型 $\Phi_{\text{ext}} \equiv \Phi_{\text{net}}$ の内容である。両者は異なる記述視点から見た同一の構造である。$\blacksquare$

\bigskip


\subsection{同値命題(誤り許容性 $\Leftrightarrow$ $\mu_\theta$)のスケッチ}


\textbf{命題}: サピエンス的知性において、誤り許容性と閾値移動能力 $\mu_\theta$ の存在は構造的に同値である。

\textbf{スケッチ($\Rightarrow$方向: $\mu_\theta$ の存在 → 誤り許容性)}:

$\mu_\theta$ は $\theta$ を動かす能力(領域組み替え (B))。$\theta$ が動くということは、現在の停止が確定的・最終的でないことを意味する。確定的でない停止は、後に組み替えられうる停止であり、したがって誤りうる停止である。よって $\mu_\theta$ の存在は誤り許容性を含意する。

\textbf{スケッチ($\Leftarrow$方向: 誤り許容性 → $\mu_\theta$ の存在)}:

誤り許容性は「停止が後に改訂されうること」を前提とする。停止の改訂は、現在の運動可能領域での停止が再検討されることを意味し、これは $\theta$ の変化(新しい停止条件の下での再評価)を要求する。$\theta$ を変化させる能力が $\mu_\theta$ である。よって誤り許容性は $\mu_\theta$ の存在を含意する。$\blacksquare$

\bigskip
